% ============================================================================
% Section 3.11 — Introduction to Scientific Notation (Virginia)
% State-specific additional topic
% ============================================================================
\section{Introduction to Scientific Notation}

\begin{stepGoal}
\begin{itemize}
    \item Write large and small numbers in scientific notation
    \item Compare numbers in scientific notation
\end{itemize}
\end{stepGoal}

\begin{stepVocabBox}
    \vocabItem{Scientific Notation}{A way to write a number as $a \times 10^n$, where $1 \leq a < 10$.}
    \vocabItem{Exponent}{The power of $10$ that tells how many places the decimal moved.}
\end{stepVocabBox}

\begin{stepsCard}[How to Write a Number in Scientific Notation]
    \stepItem{Move the decimal point so you have a number from $1$ to less than $10$.}
    \stepItem{Count how many places you moved the decimal — that is the \textbf{exponent}.}
    \stepItem{If the original number is \textbf{large} ($\geq 10$), the exponent is \textbf{positive}. If \textbf{small} ($< 1$), the exponent is \textbf{negative}.}
\end{stepsCard}

% --- Worked Example 1 ---
\begin{stepExample}{Write $7{,}200{,}000$ in scientific notation.}
    \stepShow{1} Move the decimal to get $7.2$ (between $1$ and $10$).

    \stepShow{2} The decimal moved $6$ places to the left.

    \stepShow{3} The original number is large $\to$ positive exponent: $7.2 \times 10^6$.
    \stepResult{$7.2 \times 10^6$}
\end{stepExample}

% --- Worked Example 2 ---
\begin{stepExample}{Write $0.00056$ in scientific notation.}
    \stepShow{1} Move the decimal to get $5.6$.

    \stepShow{2} The decimal moved $4$ places to the right.

    \stepShow{3} The original number is small $\to$ negative exponent: $5.6 \times 10^{-4}$.
    \stepResult{$5.6 \times 10^{-4}$}
\end{stepExample}

\mascotStepTip{Count the hops! The exponent equals how many places you moved the decimal point.}

% ============================================================================
% PRACTICE
% ============================================================================
\newpage
\begin{stepPractice}{Scientific Notation Practice}
    \resetProblems

    \stepPracticeHeader[step1Color]{Move the Decimal}
    \begin{multicols}{2}
    \prob $93{,}000{,}000 \to$ \answerBlank[2cm] $\times 10^?$
    \answer{$9.3$}
    \prob $0.0072 \to$ \answerBlank[2cm] $\times 10^?$
    \answer{$7.2$}
    \end{multicols}

    \stepPracticeHeader[step2Color]{Write in Scientific Notation}
    \begin{multicols}{2}
    \prob $93{,}000{,}000 =$ \answerBlank[3cm]
    \answer{$9.3 \times 10^7$}
    \prob $0.00056 =$ \answerBlank[3cm]
    \answer{$5.6 \times 10^{-4}$}
    \end{multicols}

    \stepPracticeHeader[step3Color]{Convert and Compare}
    \prob Convert $3.07 \times 10^5$ to standard form. \answerBlank[3cm]
    \answer{$307{,}000$}

    \wordProblem{Which is larger: $6.1 \times 10^8$ or $9.2 \times 10^7$? Explain.}{}
    \answer{$6.1 \times 10^8$ is larger because $10^8 > 10^7$.}
\end{stepPractice}
